\documentclass[11pt,a4paper]{article}

\usepackage[margin=1.1in]{geometry}
\usepackage{amsmath,amssymb,amsthm,mathtools}
\usepackage{stmaryrd}
\usepackage{tikz-cd}
\usepackage[T1]{fontenc}
\usepackage{mathpazo}
\usepackage{microtype}
\usepackage{enumitem}
\usepackage{booktabs}
\usepackage{float}
\usepackage[hidelinks]{hyperref}

\theoremstyle{definition}
\newtheorem{definition}{Definition}[section]
\newtheorem{construction}[definition]{Construction}
\newtheorem{design}[definition]{Design choice}
\newtheorem{example}[definition]{Example}
\newtheorem{notation}[definition]{Notation}

\theoremstyle{plain}
\newtheorem{proposition}[definition]{Proposition}
\newtheorem{theorem}[definition]{Theorem}
\newtheorem{lemma}[definition]{Lemma}
\newtheorem{corollary}[definition]{Corollary}
\newtheorem{conjecture}[definition]{Conjecture}

\theoremstyle{remark}
\newtheorem{remark}[definition]{Remark}
\newtheorem{problem}{Problem}

\newcommand{\Ctx}{\mathrm{Ctx}}
\newcommand{\Sim}{\mathrm{Sim}}
\newcommand{\Psh}{\mathrm{Psh}}
\newcommand{\GSLT}{\mathbf{GSLT}}
\newcommand{\iGSLT}{\mathbf{iGSLT}}
\newcommand{\ciGSLT}{\mathbf{ciGSLT}}
\newcommand{\CtxGSLT}{\mathbf{CtxGSLT}}
\newcommand{\Nerve}{\mathrm{N}}
\newcommand{\Real}{\mathrm{R}}
\newcommand{\enc}[1]{\llbracket #1 \rrbracket}
\newcommand{\code}[1]{\ulcorner #1 \urcorner}
\newcommand{\hole}{[-]}
\newcommand{\bisim}{\approx}
\newcommand{\obs}{\Downarrow}

\title{\bf Probing Universality\\[2pt]
\large Replacing the Turing Yardstick with a Relative Property\\ for Interactive GSLTs}
\author{L.\ G.\ Meredith\thanks{\texttt{lgreg.meredith@f1r3fly.io}}\\ F1R3FLY.io}
\date{July 2026}

\begin{document}
\maketitle

\begin{abstract}
Expressiveness comparisons between computational models are conventionally
anchored to a single absolute yardstick: Turing completeness. For functions this
works. For interactive systems it does not, because Turing completeness is a
statement about the input/output relation and the input/output relation is the
coarsest invariant in the linear-time/branching-time spectrum, while the
equivalences that matter for interaction sit at the other end. Calculi of identical
Turing power are routinely separated by uniform-encoding criteria, so the yardstick
is measuring something other than what is being compared.

The obvious repair --- demand behavioural rather than computational universality ---
merely relocates the problem, since the arbiter of which behavioural equivalence to
use is an adequate modal logic, and the logical apparatus available in the ambient
setting is inherited from the $\lambda$-calculus. An absolute notion is quietly
$\lambda$-relative.

This note replaces the yardstick with a relative property. One fixes a probe --- Turing
machines, $\lambda$, CCS, $\pi$, rho, or anything else --- and asks what that probe can
see of the theory under investigation. No logic is required: bisimulation is
identification in a final coalgebra, and the only primitive observation is reduction.
What must be supplied instead are conditions preventing the question from having a
trivial answer.

We recall the naive category of graph-structured lambda theories with
bisimulation-preserving morphisms and exhibit three defects: the constant map to $0$
is a morphism; the canonical encoding of the $\lambda$-calculus into $\pi$ is not a map
of terms at all, so signature preservation does not typecheck; and ambient contexts
in a richer theory observe strictly more than the encoded contexts of a poorer one.
We then rebuild the definition with contexts as the primary structure. A theory
presents a symmetric multicategory of contexts whose interfaces record sort, binding
stage, and interaction surface; transitions are labelled by contexts; and a morphism
of theories is a pseudofunctor of context multicategories, so that terms, being the
nullary contexts, are encoded by the same map that encodes contexts. Bisimulation on
the target is computed only over the image of the probe's contexts.

Two conditions are then isolated: \emph{hosting} (faithfulness), which says the
target can run the source, and \emph{exhausting} (density), which says the target is
nothing the source cannot assemble. They induce a preorder whose degrees refine
Turing completeness, the classical notion being recovered as the trace collapse ---
which above the Turing threshold has exactly one degree where the branching-time
order has many. Our main structural observation is that the two conditions are the
two halves of mutual meta-circularity: hosting yields an interpreter for the source
written in the target, exhausting an interpreter for the target written in the
source, and their two composites are meta-circular interpreters, each behaviourally
the identity and syntactically not.
\end{abstract}

\tableofcontents

\section{Introduction}
\label{sec:intro}

Ask whether one computational model is at least as expressive as another and the
answer is normally routed through a single absolute standard. Both models are
compared to the Turing machine; if both clear the bar, they are equi-expressive.
That procedure is sound for models of function computation, and it is the reason the
classification of programming languages by expressive power is so short a list.

For interactive systems it fails, and it fails for a structural reason rather than a
technical one. Turing completeness is a statement about which partial recursive
functions a model realises --- a property of its input/output relation. The
input/output relation is the coarsest invariant in the linear-time/branching-time
spectrum. The equivalences that matter for interacting systems --- bisimulation and its
refinements --- sit at the opposite end. A hypothesis stated at the trace level cannot
control a conclusion stated at the branching level. Concretely: an interpreter must
decode before it acts, decoding is silent work, and silent work relocates choice
points, so an encoding may be functionally perfect and behaviourally wrong. The gap
is not hypothetical. The separation literature exhibits pairs of calculi of identical
Turing power between which no encoding satisfying reasonable uniformity criteria
exists \cite{palamidessi,gorla,peters-glabbeek}.

The natural repair is to demand of a model not computational universality but
\emph{behavioural} universality: a self-interpreter faithful up to bisimulation. But
stated absolutely, this only relocates the difficulty. Which bisimulation? The usual
arbiter is adequacy of a modal logic, and the logical constructions available in the
ambient setting are inherited from the $\lambda$-calculus. An absolute behavioural
notion is quietly $\lambda$-relative, which is precisely the parochialism one was
trying to escape.

This note therefore replaces the yardstick rather than recalibrating it. One fixes a
\emph{probe} --- Turing machines, $\lambda$, CCS, $\pi$, rho, or anything else --- and
asks what that probe can see of the theory under investigation. Universality becomes
a relation, not a predicate. No logic is required: bisimulation is identification in
a final coalgebra, and the only primitive observation is reduction. What must be
supplied instead are conditions preventing the question from having a trivial answer,
and most of this note is about finding the right ones. They turn out to be two, they
are standard categorical properties rather than a checklist, and together they say
precisely that source and target can write meta-circular interpreters for one
another.

The classical notion is not discarded. It reappears in Section~\ref{sec:tc} as the
trace collapse of the relative order --- a single degree sitting above a threshold,
where the branching-time order has many. Turing completeness is thus the shadow of
the relative property, and the separations that the classical notion cannot see are
exactly the ones the shadow loses.

\paragraph{Outline.}
Section~\ref{sec:naive} fixes the class of theories through a three-rung ladder,
spells out the shape of the core rewrite rule, shows why the machine models require
impedance matching to enter the category at all, and explains via the McBride derivative
why equations on the interaction constructor force surfaces to be reified --- which
matters because the contexts whose derivative is at issue are exactly the labels of
the transition systems of Section~\ref{sec:contexts}. It closes with the defects of
the naive notion of morphism.
Section~\ref{sec:contexts} rebuilds the category with contexts as the primary
structure. Section~\ref{sec:conditions} isolates the two conditions.
Section~\ref{sec:tc} relates the resulting order to Turing completeness.
Section~\ref{sec:meta} proves the two-interpreter theorem. Section~\ref{sec:open}
lists what is not yet settled.

\section{The naive setting and its defects}
\label{sec:naive}

\subsection{A ladder of theories}

The theories in play are pinned down in three stages, each rung adding structure
and each still accommodating the usual suspects. The progression matters for what
follows: the constructions of Sections~\ref{sec:contexts}--\ref{sec:conditions}
need only the second rung, while the interpreters of Section~\ref{sec:meta} need
the third, and knowing which is which is what keeps the hypotheses honest.

\paragraph{Rung one: graph-structured lambda theories.}

\begin{definition}[GSLT]
\label{def:gslt}
A \emph{graph-structured lambda theory} $G$ consists of a sorted signature with
binding; a set of terms generated by that signature; a structural congruence
$\equiv$ making the term set a graph-like quotient, so that terms are considered up
to the commutative-monoid laws of parallel composition, scope extrusion, and the
like; and a rewrite relation $\to$ on the quotient.
\end{definition}

At this rung almost everything qualifies --- $\lambda$, CCS, $\pi$, rho, the ambient
calculus, and Turing machines with $\equiv$ trivial. That breadth is the problem:
Definition~\ref{def:gslt} constrains the \emph{shape} of a theory's dynamics not at
all, so nothing can be said uniformly about how one theory simulates another.

\paragraph{Rung two: interactive.}

The second rung asks only for a \emph{site of interaction}. Following
\cite[Def.~2.1]{cost}:

\begin{definition}[Interactive GSLT]
\label{def:igslt}
$G$ is \emph{interactive} if it carries, over and above the data of
Definition~\ref{def:gslt}, a distinguished interacting sort $P$ of programs; a binary
\emph{interaction constructor} $K(P,E)$ representing the interaction of a program with
an environment $E$, an operand of the interacting sort, with $E = P$ in the
process-calculus instances; and at least one base rewrite rule whose left-hand side
features $K$, generating a labelled transition system on $P$ modulo $\equiv$ on which
bisimulation is the intended equivalence.
\end{definition}

For the $\lambda$-calculus $K$ is application; for CCS, $\pi$, and rho it is parallel
composition; for the interaction categories it is composition along a shared
interface. Nothing about surfaces or continuations is asserted at this rung. What a
bare GSLT lacks and an interactive one supplies is exactly a named site of
interaction and a rule that fires there.

\begin{remark}[Where the machine models fail]
\label{rem:machines}
The requirement that $E$ be an operand of the interacting sort is not a
technicality. A Turing machine does have an intuitive cut --- between the automaton
and the tape --- but automaton and tape are of different kinds, so there is no binary
$K$ on the interacting sort of which the transition rule is an instance. Turing
machines therefore admit no interaction-cut presentation naively, and neither do
Mealy or Moore machines, whose cut is between a transducer and a stream. They are
GSLTs; they are not interactive ones. Section~\ref{sec:tc} takes this seriously: the
classical yardstick is not merely a coarse probe, it is not natively an object of the
category at all, and must be \emph{presented} as one --- by a choice, such as encoding
the tape as a process, that is in no way canonical.
\end{remark}

\paragraph{Rung three: continued.}

The third rung decomposes both sides of the cut. It is here, and not before, that
surfaces and continuations appear.

\begin{definition}[The interaction cut]
\label{def:cut}
An interactive GSLT is presented in \emph{interaction-cut form} when its generating
rule is written
\[
K\bigl(\,K_p(I,C_p),\; K_e(J,C_e)\,\bigr) \;\longrightarrow\; K'\bigl(\,C_p',\,C_e'\,\bigr),
\qquad
(C_p',C_e') = \mathrm{compute}(I,J,C_p,C_e)
\tag{$\dagger$}
\]
naming the interaction constructor $K$; two introductions $K_p(I,C_p)$ and
$K_e(J,C_e)$, each separating an \emph{interaction surface} ($I$, $J$ --- the part that
must match for the cut to fire) from a \emph{continuation} ($C_p$ the program
continuation, $C_e$ the environment continuation); the partial \emph{contraction}
$\mathrm{compute}$, which on a surface match combines the continuations into
residuals; and the constructor $K'$ repackaging them.
\end{definition}

$\mathrm{compute}$ is Milner's pseudo-application: applied to $\lambda y.P$ and
$[u,Q]$ it returns $P\{u/y\} \mid Q$, factoring into the substitution by which the
datum migrates from the environment continuation into the program continuation
through the binder, and the release of the environment's residual. Binding is one
mode of migration, not a feature of the cut as such; CCS realises the shape with
pseudo-application degenerating to pure release.

\begin{definition}[Surfaces and nearness]
\label{def:surface}
Surfaces inhabit a sort $\mathrm{Surf}$. When $K$ is non-free the match condition is
carried by a \emph{nearness operator}
\[
\mathrm{near}(I,J),
\]
partial, and returning --- when defined --- the surface at which $I$ and $J$ meet.
Interaction is gated on its being defined.
\end{definition}

Nearness is thus \emph{surface-valued} rather than boolean: two introductions do not
merely match or fail to match, they produce a surface, and that surface locates the
resources the cut may draw on. It specialises across the spectrum as complementarity
of actions in CCS, name-equality of subjects in $\pi$ and rho, and interface
agreement in the interaction categories.

\begin{definition}[Continued interactive GSLT]
\label{def:cigslt}
An interactive GSLT is \emph{continued} when equipped with
\begin{enumerate}[label=(\roman*),leftmargin=2.4em]
\item a presentation of its dynamics in interaction-cut form $(\dagger)$, naming
$K$, $K_p$, $K_e$, $K'$ and the partial contraction $\mathrm{compute}$; and
\item a \emph{section}: a computable canonical-form function on the interacting sort,
i.e.\ a computable section of the $\equiv$-quotient.
\end{enumerate}
$\ciGSLT$ is the category of these, with morphisms the theory maps that are
bisimulation-preserving on the interacting sort.\footnote{The definition of a
continued interactive GSLT in \cite[Def.~6.1]{cost} carries a third clause, required
there for a construction not treated in this note. Nothing below depends on it, and
we omit it.}
\end{definition}

The adjective names a strict enrichment, not a restriction: (i) and (ii) are data
added on top of an interactive GSLT, and forgetting them recovers it. It records that
the theory carries the structure a construction \emph{over} it will need, not that
any such construction has been performed.

\begin{remark}[Two carriers of the surface, and why $\lambda$ survives]
\label{rem:regimes}
Definition~\ref{def:cigslt} appears to exclude the $\lambda$-calculus, which has no
sort of surfaces and no nearness operator. It does not, because a surface may be
carried in either of two ways.
\begin{enumerate}[label=(\alph*),leftmargin=2.0em]
\item \emph{Structurally, when $K$ is free.} Position is then unforgeable: nothing
drifts into contact with what it does not already neighbour. The surface is not
recorded by any explicit $I$ or $J$ but by the structural context $K([\,],S)$, and
$\mathrm{near}$ is trivially defined --- being in the hole is being near. $\lambda$ is
exactly this case on the environment side: $\mathrm{App}$ is free, and its rigidity
supplies the surface that no explicit $J$ records. A nominally absent surface is not
a missing match condition but one carried by the geometry of $K$.
\item \emph{Nominally, when $K$ is non-free.} Any equation on $K$ loosens position:
associativity re-brackets adjacency, a unit law inserts or deletes neighbours, and
associative--com\-mu\-ta\-tiv\-i\-ty is the limit at which no positional structure
survives at all --- the freely mixed bag, where rho with $K = {\mid}$ lives. Adjacency having
become cheap and non-exclusive, taking proximity as authority is ambient authority
and is a leak. The surface must therefore be reintroduced nominally, supplied by
$\mathrm{near}$.
\end{enumerate}
The two carriers are one notion seen twice; only the carrier changes, from the
geometry of $K$ to an explicit operator. A free $K$ and a freely mixed $K$ are the
poles, and every $K$ carrying equations sits on the nominal side between them.
\end{remark}

\begin{table}[h]
\centering
\small
\begin{tabular}{@{}lllll@{}}
\toprule
Instance & Surfaces $I,J$ & Prog.\ cont.\ $C_p$ & Env.\ cont.\ $C_e$ & Contraction $\mathrm{compute}$ \\
\midrule
CCS
  & $a,\bar a$ (compl.)
  & $P$ & $Q$
  & pure release $P \mid Q$ \\
rho / $\pi$ (sync.)
  & $x,x$ (name-eq.)
  & $\lambda y.P$ & $[Q_1,Q_2]$
  & $P\{Q_1/y\} \mid Q_2$ \\
rho / $\pi$ (async.)
  & $x,x$ (name-eq.)
  & $\lambda y.P$ & $[Q,-]$
  & $P\{Q/y\}$ \\
$\lambda$
  & $x,-$ ($K_e$ degen.)
  & $\lambda x.M$ & $N$
  & $M\{N/x\}$ \\
Ambients
  & $n$ (open $n$ vs.\ $n[\cdot]$)
  & $P$ & $Q$
  & boundary dissolve / move \\
Interaction cats.
  & shared interface
  & process & process
  & interface composition \\
\bottomrule
\end{tabular}
\caption{Anchoring instances across the migration spectrum, after
\cite[Tab.~1]{cost}. A surface entry of $-$ denotes a nominally degenerate surface,
carried structurally by the rigidity of $K$ rather than absent; an
environment-continuation entry of $-$ denotes a degenerate continuation, which is
what distinguishes asynchronous from synchronous output. Four migration modes
appear: none (CCS, interaction categories), binding (rho, $\pi$, $\lambda$), and
spatial restructuring (ambients), where the contraction relocates or dissolves a
boundary rather than substituting.}
\label{tab:instances}
\end{table}

The instances differ in how much information migrates between environment and program
under contraction, and this migration spectrum is what makes the class wide rather
than narrow. CCS realises the shape with no migration at all and is in that sense the
paradigmatic instance; $\pi$ and rho realise it with migration by binding;
$\beta$-reduction with migration and no environment residual --- the environment
introduction has a degenerate surface and its continuation carries no residual, so the
argument is its own continuation and pseudo-application is ordinary application.

\begin{remark}[Which clause this note uses]
\label{rem:whichhalf}
Sections~\ref{sec:contexts} and \ref{sec:conditions} need rung two, together with
clause (i) only insofar as surfaces sharpen the interfaces of
Definition~\ref{def:interface}; a reader who drops the surface component throughout
loses resolution in the labels and nothing else. Clause (ii), the section, is what
Section~\ref{sec:meta} needs: it is what makes the interpreters exist as terms rather
than as relations.
\end{remark}

\begin{proposition}[The section makes $\Ctx(G)$ effectively presented]
\label{prop:effctx}
If $G$ is continued then $\Ctx(G)$ admits codes: equality of contexts is decidable on
representatives, and hole-filling is computable.
\end{proposition}

\begin{proof}
Clause (ii) of Definition~\ref{def:cigslt} picks a canonical representative of each
$\equiv$-class, so a context is presented by a representative term with labelled
holes. Hole-filling is a syntactic substitution on representatives followed by
re-application of the section. Stage and surface data
(Definition~\ref{def:interface}) are finite and carried along.
\end{proof}

Proposition~\ref{prop:effctx} is why clause (ii) cannot be dispensed with. Without
it, $\Ctx(G)$ still exists as a mathematical object --- quotient multicategories are
unproblematic --- but nothing about it is nameable, and the interpreters of
Section~\ref{sec:meta} would exist only as relations between terms rather than as
terms of the calculi themselves.

\paragraph{The usual suspects.}
Each rung retains the examples that matter, and loses exactly the machine models.

\begin{center}
\small
\begin{tabular}{@{}lcccp{4.2cm}@{}}
\toprule
Theory & GSLT & Interactive & Continued & $K$, and the section \\
\midrule
Turing machines & \checkmark & --- & --- & cut is automaton/tape: different kinds \\
Mealy, Moore & \checkmark & --- & --- & cut is transducer/stream: different kinds \\
$\lambda$ ($\beta$) & \checkmark & \checkmark & \checkmark & $\mathrm{App}$; de Bruijn form, no channel constructor needed \cite[\S5.2]{cost} \\
CCS & \checkmark & \checkmark & \checkmark & $\mid$; monoid normal form \\
$\pi$ & \checkmark & \checkmark & \checkmark & $\mid$; as CCS plus $\alpha$-normalisation \\
rho & \checkmark & \checkmark & \checkmark & $\mid$; $\#$ fuses section and quotient \\
Interaction cats. & \checkmark & \checkmark & \checkmark & composition along an interface \\
\bottomrule
\end{tabular}
\end{center}

The two exclusions are the point of the rung, not an embarrassment of it. What the
machine models lack is not computational power but a cut between operands of the
interacting sort, and it is precisely that absence which makes them unable to say
anything about how one interacting system simulates another. That $\lambda$ passes
clause (ii) without possessing anything like rho's $@\cdot$ is the sharp test of the
same point from the other side: what is required is a section, not a channel
constructor \cite[\S5.2]{cost}.

\begin{remark}[The ladder is not the probe order]
\label{rem:ladder-vs-order}
Ascending the ladder narrows the class of theories under discussion; it says nothing
about their relative expressiveness. $\lambda$ and rho are both continued, and
Section~\ref{sec:tc} will place them at very different degrees. The ladder fixes the
category; the probe order fixes the ordering within it.
\end{remark}

Throughout, $\obs$ denotes a fixed primitive observation on terms. We take
$P \obs$ to mean that $P$ has a reduction, or --- when a success predicate is
available --- that $P$ reports success. This is the only observation we assume, and
it is intrinsic to the rewrite relation. In particular no modal logic is used
anywhere in this note.

\begin{definition}[Naive morphism]
\label{def:naive}
A \emph{naive morphism} $f : G \to H$ is a function on terms satisfying
$P \bisim_G Q \implies f(P) \bisim_H f(Q)$, where $\bisim$ denotes the intended
weak bisimulation of each theory. Write $\GSLT_{\mathrm{naive}}$ for the resulting
category.
\end{definition}

This is the definition implicitly in play when one says that morphisms of $\ciGSLT$
are ``bisimulation-preserving and quote-faithful''. Quote-faithfulness is an extra
condition specific to reflective theories and does not repair what follows.

\subsection{Impedance matching: why the machine models fail}
\label{sec:impedance}

Remark~\ref{rem:machines} asserted that Turing, Mealy, and Moore machines are GSLTs
but not interactive ones. The claim is worth working, both because the classical
yardstick is built from them and because the repair --- which is always available ---
turns out to introduce exactly the pathology this note exists to diagnose.

\paragraph{They are GSLTs.}
A Turing machine presents as a GSLT without difficulty. Take sorts for control
states, tape alphabet, and configurations; terms of the configuration sort are
triples $\langle q, w, i\rangle$ of state, tape contents, and head position;
$\equiv$ is trivial; and $\to$ is the transition relation. Everything
Definition~\ref{def:gslt} asks for is present. Mealy and Moore machines present
similarly, with configurations pairing a state with an input stream.

\paragraph{They are not interactive.}
Definition~\ref{def:igslt} asks for a binary $K(P,E)$ whose operands both inhabit the
interacting sort, and for a rule headed by $K$. The natural cut in a Turing machine
is between the automaton and the tape --- but the automaton inhabits the sort of
control states and the tape inhabits the sort of words, so any $K$ pairing them is a
constructor of a heterogeneous product, not a binary operation on one sort. The
asymmetry is of \emph{kind}, not merely of role: in an interaction cut program and
environment are interchangeable in kind while distinguished in role, whereas a tape
can never be an automaton. Consequently the transition relation is not a $K$-headed
rule closed under $K$-contexts, and the labelled transition system it generates does
not live on an interacting sort.

Mealy and Moore machines fail for the same reason and for a second one besides.
Their dynamics is a function on state and input rather than a rewrite of a term;
there is no term algebra on which the step is congruent, so even granting a cut there
is nothing for it to be closed under. The failure is a failure of compositionality
before it is a failure of typing.

\paragraph{Impedance matching, and what it costs.}
The repair is standard and always available: reflect the environment into the
interacting sort. Encode the tape as a process --- a chain of cells communicating with
a head process --- and both operands of $K = {\mid}$ become processes. Call the
result a \emph{process presentation}, $\mathrm{TM}^\sharp$.

Three observations, and the third is the one that matters here.

\begin{enumerate}[leftmargin=1.6em]
\item The encoding is a choice. Cells may be a linked list, an array of indexed
agents, or a pair of stacks; each yields a different $\Ctx(\mathrm{TM}^\sharp)$ and
hence a different probe. Nothing internal to the machine model selects among them.
\item The encoding is not conservative on steps. One machine transition becomes
several cut firings --- request, reply, commit --- so the induced transition system is
not the machine's but a refinement of it.
\item Those extra firings are administrative, and administrative steps relocate
choice points. That is precisely the mechanism by which a functionally correct
interpreter becomes behaviourally wrong (Section~\ref{sec:intro}).
\end{enumerate}

So the admission ticket that lets the classical yardstick into the category
introduces the very pathology the yardstick is constitutionally unable to see. This
is not an argument against the machine models; it is an argument against using them
as the standard against which interacting systems are measured.

\subsection{Locations as labels: the derivative and the reification of surfaces}
\label{sec:derivative}

\paragraph{Why this subsection is here.}
Section~\ref{sec:clts} will label transitions by \emph{contexts}: a term's behaviour
is recorded not by abstract actions drawn from some separately postulated alphabet,
but by the contexts that can make it move. A context is a location. So the labels of
every transition system in this note are locations, and a labelling is informative
only to the extent that its labels are distinguishable. The question of when a
location is a genuine coordinate on a term is therefore the question of how much
resolution the probe order of Sections~\ref{sec:conditions}--\ref{sec:tc} has at its
disposal, and it is settled here, before any of that machinery is built. The chain is
short and worth stating in advance:
\[
\begin{array}{rcl}
\text{contexts separate positions} &\iff& \mathrm{plug} \text{ is injective}\\[2pt]
&\Longrightarrow& \text{context labels carry information}\\[2pt]
&\Longrightarrow& \text{the induced bisimulation is fine.}
\end{array}
\]
Everything below is an analysis of when the first line holds, what happens when it
fails, and what must be added to the theory to compensate.

Remark~\ref{rem:regimes} said that a free $K$ carries the surface structurally, as
the context $K([\,],S)$, and that equations on $K$ force the surface to be carried
nominally instead. The erosion can now be stated exactly, and the right instrument is
McBride's derivative \cite{mcbride}.

\paragraph{Location is a zipper, and a zipper is not a term.}
Let the terms be $T = \mu X.\,F(X)$ for a signature functor $F$. The one-hole
contexts are not $\partial F$ but paths of layers: writing $\partial F(T)$ for a
single layer with one slot vacated and the rest filled with terms,
\[
\Ctx_T \;=\; \mathrm{List}\bigl(\partial F(T)\bigr),
\qquad
\mathrm{plug} : \Ctx_T \times T \to T .
\]
This is Huet's zipper \cite{huet} read algebraically via McBride's derivative
\cite{mcbride}. Two points must be kept apart, and running them together is the error
this subsection is written to avoid. First, $\Ctx_T$ and $T$ are \emph{different
types}. Second, a location is a coordinate on a term exactly when $\mathrm{plug}$ is
injective --- distinct contexts must name distinct holes --- and that is a statement
about $\mathrm{plug}$, not about any isomorphism between $\Ctx_T$ and $T$.

Layer derivatives for the constructors at issue. The right-hand column says what the
derivative \emph{is}; it does not say that any of these is a surface.
\[
\begin{array}{llll}
\text{prefix} & F_0(X) = A \times X & \partial F_0(X) = A & \text{an action}\\
\text{restriction} & F_0(X) = \mathcal{N} \times X & \partial F_0(X) = \mathcal{N} & \text{a name}\\
\text{abstraction} & F_0(X) = \mathrm{Var} \times X & \partial F_0(X) = \mathrm{Var} & \text{a variable}\\
\text{free binary }K & F_0(X) = X \times X & \partial F_0(X) = X + X & \text{a sibling, tagged}\\
\text{AC }K & F_0(X) = \mathrm{Bag}\,X & \partial F_0(X) = \mathrm{Bag}\,X & \text{the rest of the bag}
\end{array}
\]
Only the two $K$-rows bear on surfaces at all, since only $K$ is the site of the cut.
The abstraction row is listed because it is what makes erasure by binding possible
below, and because the variable it produces is easily and wrongly taken for a
surface.

\paragraph{Free $K$: the derivative is a coordinate system.}
For a free binary $K$ the layer derivative is $T + T$ --- a sibling subterm together
with a tag saying which side the hole is on --- and it is the \emph{tag}, not the
sibling, that carries the positional information. $\mathrm{plug}$ is injective, so
$\Ctx_T$ is a coordinate system on terms, and $K([\,],S)$ is a context that names its
position uniquely. This is the precise content of position being unforgeable: the
surface is the location, and no reification is needed because $\Ctx_T$ is already
doing the work.

\begin{remark}[The surface of a $\beta$-redex, and what the variable is instead]
\label{rem:lambdasurface}
The $\lambda$-calculus is the case to state carefully, since it is the one where the
surface is easiest to misidentify. $K = \mathrm{App}$ is free, so by the paragraph
above the surface is supplied by the location. An application comprises a location
together with a function and an argument; the \emph{location is the surface}, the
function is the program and the argument is the environment --- or, dually, the
argument is the program and the function is the environment. The location itself is
symmetric; which operand is assigned which role is an orientation, not something the
syntax fixes. What is not available in either orientation is a nominal surface,
precisely because $\mathrm{App}$ is rigid.

A bound variable is therefore \emph{not} a surface. It is nonetheless closely
related, and the relation is worth naming, because both are derivative data. The
surface is the location at which program and environment come into contact --- the
$\partial$ of the $K$-layer. The variable is the name of the locations \emph{inside}
the program continuation at which the environment's contribution will be delivered:
the occurrences of $x$ in $M$ are a multiset of holes in $M$, and $\lambda x.M$ is
exactly that multi-hole context internalised as a term with $x$ naming its holes.
Contact site and delivery site are both locations, at different levels of the term
and playing different roles in the cut.

This also sharpens Table~\ref{tab:instances}, whose $\lambda$ row records the
program-side surface as $x$ following \cite[Tab.~1]{cost}. Read through the rigidity
analysis, that entry is naming the delivery site rather than the contact surface; the
contact surface on \emph{both} sides is the $\mathrm{App}$ location, which is why the
environment-side entry is nominally degenerate rather than absent.
\end{remark}

\paragraph{Equations erode injectivity.}
Once $K$ carries equations the derivative is computed in the quotient, and what
degrades is injectivity of $\mathrm{plug}$ --- not the type of contexts.
\begin{itemize}[leftmargin=1.4em]
\item \emph{Associativity.} Trees collapse to lists and formerly distinct bracketings
of the same neighbourhood become the same context.
\item \emph{Units.} Neighbours may be inserted or deleted, so one location is
reachable from strictly more decompositions.
\item \emph{Associativity with commutativity.} Lists collapse to bags and
$\partial(\mathrm{Bag}\,T) \cong \mathrm{Bag}\,T$: the derivative has forgotten
\emph{which} occurrence was vacated. For a bag containing $n$ copies of a subterm
there are $n$ decompositions no context distinguishes, and $\mathrm{plug}$ is
radically non-injective.
\end{itemize}
At the AC pole the zipper has ceased to be a coordinate system. This is what
``adjacency has become cheap and non-exclusive'' amounts to formally: not that
positions are hard to tell apart, but that $\Ctx_T$ no longer separates them.

Read through the chain above, the consequence is not merely that authority leaks. It
is that the labels of the transition system collapse. Two contexts that ought to
witness different behaviour become the same label, so the induced bisimulation
identifies terms it should distinguish, and every probe built from the theory
inherits the loss. Equations on $K$ are thus not a local nuisance about capability;
they are a global reduction in the resolution of everything the theory can observe or
be observed by.

\paragraph{Two erasures, only one of them benign.}
It is tempting, having computed $\partial(\mathrm{Bag}\,T) \cong \mathrm{Bag}\,T$, to
say that context and term have become the same thing. They have not, and the
temptation conceals a genuine distinction between two ways the type barrier can come
down.

\emph{Erasure by binding.} In a calculus with a binder and a substitution, a one-hole
context is representable \emph{as a term}: $C[\,] \mapsto \lambda x.C[x]$, with
$\mathrm{plug}$ realised by substitution. Here $\Ctx_T$ genuinely embeds into $T$ (or
into an abstraction sort over $T$), the embedding is faithful, and nothing is lost.
This is available in $\lambda$, $\pi$, and rho, and in rho it is available twice over,
since $@$ makes the representing term quotable in turn. Note that the binder is here
reifying the \emph{delivery} locations inside a continuation, which is dual to the
reification of the \emph{contact} location forced at the AC pole
(Remark~\ref{rem:lambdasurface}); the two reifications have different targets and
should not be run together. Note what this is: the ability
of a theory to internalise its own contexts, which is exactly condition (I) of
Section~\ref{sec:tc} seen from the syntax side. The erasure is not a convenience; it
is a structural property that some theories have and others do not.

\emph{Erasure by collapse.} At the AC pole the layer derivative is term-shaped for an
entirely different reason: a bag with a hole is just a smaller bag, because the hole
carries no information. Nothing has been internalised; something has been forgotten.
The two situations look alike on the page and are opposite in character --- one is a
faithful embedding, the other a quotient.

\paragraph{CCS: the distinction is irreducible.}
CCS has no binder in the substitutive sense. Restriction hides a name but is never
instantiated, and prefix binds nothing, so there is no abstraction to form and no
substitution to plug with; erasure by binding is unavailable. And the layer
derivatives above show why erasure by collapse does not cover the gap: the prefix
layer differentiates to an \emph{action} and the restriction layer to a \emph{name},
neither of which is a term. A CCS context is a path made of actions, names, and
residual bags, and that is irreducibly a different type from a CCS process.

This matters for how the present section may be used. The paradigmatic instance of
the interaction cut is the one calculus in which the context/term distinction cannot
be erased, so no argument in this note may presuppose the erasure. Accordingly the
erosion criterion above is stated purely as failure of injectivity of
$\mathrm{plug}$, which is meaningful whether or not $\Ctx_T$ embeds into $T$.

\paragraph{Hence reification.}
When $\mathrm{plug}$ ceases to be injective the surface can no longer be read off the
context, and must be \emph{reified} as an explicit datum $I : \mathrm{Surf}$ carried
by the introduction, with $\mathrm{near}(I,J)$ supplying the match condition the
zipper can no longer supply. Reification is forced, and forced exactly at the point
injectivity fails. What reification restores is a coordinate system indexed by the
reified surface in place of the collapsed one --- and with it, labels that once again
separate. The nominal surface is, in this reading, precisely the information the
derivative could no longer carry, put back by hand so that context labels retain
content.

\begin{remark}[Consequences downstream]
\label{rem:derivconseq}
Three things in the sequel follow from this subsection rather than being independent
stipulations. First, $\Ctx(G)$ of Definition~\ref{def:ctx} is a multicategory of
zippers and its members are the transition labels of Section~\ref{sec:clts}, so when
$\mathrm{plug}$ loses injectivity the label set shrinks and the induced bisimulation
coarsens; this is Remark~\ref{rem:Kequations}, now derived rather than observed. Second, interfaces
carry a surface component (Definition~\ref{def:interface}) for exactly this reason:
if the zipper separated positions, sort and stage would suffice. Third, and least
expected, erasure by binding is the syntactic face of internal representability:
a theory can code its own contexts precisely when it can abstract over a hole, which
is why $\lambda$, $\pi$, and rho satisfy condition (I) of
Definition~\ref{def:I} and CCS does not. The separation of
Corollary~\ref{cor:pinotccs} and the failure of erasure in CCS are one fact.
\end{remark}

\subsection{Defect one: preservation is the cheap direction}

\begin{proposition}
\label{prop:trivial}
For any $G$ and any $H$ possessing an inert term $0$ with $0 \not\obs$, the constant
map $P \mapsto 0$ is a naive morphism.
\end{proposition}

\begin{proof}
Immediate: the hypothesis of the implication is discharged by a constant conclusion.
\end{proof}

The trivial encoding preserves bisimulation because \emph{preservation} is the cheap
direction. What one wants is \emph{reflection},
$f(P) \bisim_H f(Q) \implies P \bisim_G Q$, i.e.\ faithfulness on the bisimulation
quotient. This kills the constant map at once. It does not, however, suffice.

\begin{remark}[Table lookup]
\label{rem:lookup}
If $H$ is Turing complete, enumerate the terms of $G$ and define $f$ by case
analysis on the G\"odel index, choosing pairwise non-bisimilar images. The result is
faithful and worthless: all the content sits in the metatheoretic map, none in $H$.
Faithfulness must therefore be paired with a condition forbidding the encoding to
inspect its argument. The obvious such condition --- that $f$ be a homomorphism of
the signature --- is refuted by the next subsection.
\end{remark}

\subsection{Defect two: signature preservation does not typecheck}

Consider Milner's encoding of the $\lambda$-calculus into $\pi$ \cite{milner-fap}.
In the call-by-name presentation,
\[
\enc{\lambda x.M}p = p(x).p(q).\enc{M}q,
\qquad
\enc{x}p = \overline{x}\langle p\rangle,
\]
\[
\enc{MN}p = (\nu q)\bigl(\enc{M}q \mid (\nu y)(\overline{q}\langle y\rangle.\overline{q}\langle p\rangle \mid {!}\,y(r).\enc{N}r)\bigr).
\]

Two things are visible.

First, $\enc{-}$ is not a function from terms to terms. Its target is
name-abstracted processes: $\enc{-} : \Lambda \to (\mathcal{N} \multimap \Pi)$,
each term being encoded relative to a continuation name at which its value is
delivered. A signature homomorphism $\Lambda \to \Pi$ is not merely false here; it
is not a well-formed demand.

Second, even granting the parameterisation, no operator of $\Pi$ corresponds to
application. The application clause introduces restriction, replication, and a
two-step forwarding protocol. What is uniform about the clause is that a fixed
\emph{context} is used, with the encodings of the subterms plugged into its holes ---
a context-level, not a signature-level, regularity.

\begin{remark}[The kernel]
\label{rem:kernel}
Nor is the encoding faithful. The equivalence it induces on $\lambda$-terms is,
by Sangiorgi's analysis of the lazy encoding \cite{sangiorgi-lazy}, L\'evy--Longo
tree equality; the call-by-value variant induces an eager-normal-form equivalence.
Both are strictly coarser than $\beta$-conversion. So the canonical example of a
good encoding fails both of the obvious repairs: it is not a signature homomorphism,
and it is not faithful.
\end{remark}

The moral we draw is the one that motivates the whole of Section~\ref{sec:contexts}.
The point of moving to bisimulation is to break with syntactic constraints; syntactic
constraints should therefore re-enter only as optional refinements, never as the
primitive. What survives the break is not the signature but the \emph{theory of
contexts}.

\subsection{Defect three: ambient contexts observe too much}

Suppose $\iota : \pi \to \text{rho}$ is any reasonable encoding. Rho possesses a
reflection operator $@$ turning a process into a name, and contexts built from it
can inspect the code of their hole in ways no $\pi$ context can. Consequently there
are $\pi$-terms $P \bisim_\pi Q$ whose images are separated by some rho context. If
bisimulation on the target is computed over \emph{all} rho contexts, no encoding of
$\pi$ into rho preserves bisimulation, and the naive category records the fact that
rho is more expressive than $\pi$ as the \emph{absence} of any morphism between them.
That is exactly backwards: the extra expressive power of the target should be
recorded as data about the encoding, not as an obstruction to its existence.

\begin{remark}[Two dual failures]
\label{rem:dual}
Defects two and three are dual and must not be conflated.
\begin{itemize}[leftmargin=1.4em]
\item \emph{Under-discrimination}: the encoding identifies source terms the source
distinguishes. Milner's $\lambda \to \pi$, whose kernel is L\'evy--Longo equality.
\item \emph{Over-discrimination}: the target's own contexts separate images the
source considers equal. $\pi \to \text{rho}$.
\end{itemize}
Restricting the class of observing contexts can only make the target equivalence
\emph{coarser}. It therefore repairs over-discrimination completely and
under-discrimination never. Two different instruments are required, and
Section~\ref{sec:conditions} supplies both: relativisation of the label class for the
first, and an image factorisation recording the kernel for the second.
\end{remark}

\section{Contexts as the primary structure}
\label{sec:contexts}

\subsection{Stages}

A first attempt at the objects of the context multicategory would take them to be
lists of sorts together with the name-arity data recording which free names a term
may use. That is the wrong definition, and wrong in the way this note exists to
avoid: it presupposes a sort of names with substitution behaviour, which is
structure imported from $\pi$ and rho rather than structure common to all
interactive GSLTs. It is ill-formed for the $\lambda$-calculus, where binders bind
variables and not channels; it is vacuous for Turing machines; and it is inert for
CCS, whose restriction operator hides a name but is never instantiated, so that
there is no substitution for the arity data to track. The correct definition is
parametric in the binding signature and degenerates, rather than breaks, in each of
these cases.

\begin{definition}[Stage category]
\label{def:stage}
The binding signature of $G$ determines a small category $\mathbb{F}_G$ of
\emph{stages}: objects are the finite name or variable contexts in which a term of
$G$ may be formed, and morphisms are the admissible maps between them. We require
that passing beneath a binder be witnessed by an inclusion
$n \hookrightarrow n \sqcup \{a\}$, and permit $\mathbb{F}_G$ to be terminal.
\end{definition}

\begin{remark}[Stages across the instances]
\label{rem:stages}
The stage category, not a clause in the definition of interface, is where the
instances differ.

\begin{center}
\begin{tabular}{@{}llp{5.0cm}@{}}
\toprule
Theory & $\mathbb{F}_G$ & Action \\
\midrule
TM, finite CCS & terminal & none \\
CCS with restriction & $\mathbb{I}$: finite sets, injections & renaming; binders hide but never instantiate \\
$\lambda$ & free-variable contexts & substitution of terms for variables \\
$\pi$ & $\mathbb{F}$: finite sets, all maps & substitution; names may be identified \\
rho & fixed point with $\mathcal{N} \cong @\,\mathrm{Proc}$ & reflection; the name object is the term object \\
\bottomrule
\end{tabular}
\end{center}

This is already an expressiveness stratification, and Section~\ref{sec:tc} will
show that some of the separations in the probe order are visible at the level of
$\mathbb{F}_G$ alone, without descending to particular terms.
\end{remark}

\subsection{The context multicategory}

\begin{definition}[Interface]
\label{def:interface}
An \emph{interface} of $G$ is a triple $(s,n,\sigma)$ with $s \in S$ a sort, $n$ an
object of $\mathbb{F}_G$, and $\sigma$ a finite multiset of elements of
$\mathrm{Srf}$ --- the surfaces offered at that position.
\end{definition}

\begin{remark}[Why interfaces carry surfaces]
\label{rem:surfaceinterface}
A context can interact with its hole exactly when it offers a surface near, in the
sense of Definition~\ref{def:surface}, to one the hole exposes. Recording $\sigma$
is therefore not bookkeeping: it is what makes hole-filling
\emph{interaction-relevant} rather than merely syntactic. Composition of contexts
composes surface multisets, with pairs consumed by $\smile$ replaced by their
surface of contact. Sort says what kind of thing sits at the position, stage says
which names are in scope there, and surface says what interaction is on offer;
Section~\ref{sec:tc} separates theories using the second, and the third is what the
labels of Section~\ref{sec:clts} ultimately record.
\end{remark}

\begin{definition}[Context multicategory]
\label{def:ctx}
The \emph{context multicategory} $\Ctx(G)$ has interfaces as objects. A morphism
\[
C \;:\; (s_1,n_1,\sigma_1),\dots,(s_k,n_k,\sigma_k) \longrightarrow (s,n,\sigma)
\]
is a term of sort $s$ formed at stage $n$ offering surfaces $\sigma$, with $k$
labelled holes taken modulo $\equiv$, where hole $i$ has sort $s_i$, occurs beneath a
chain of binders exhibiting a map $n \to n_i$ in $\mathbb{F}_G$, and is offered the
surfaces $\sigma_i$ by the surrounding term. Composition is hole-filling, defined
when the stage and surface data of the filler agree with those recorded at the hole;
the identity on $(s,n,\sigma)$ is the bare hole $\hole$; the symmetric group acts by
permuting hole labels.
\end{definition}

\begin{remark}[Capture is data, not a defect]
\label{rem:capture}
Hole-filling is deliberately \emph{not} capture-avoiding: $(\nu a)\hole$ captures,
and it must, since that is what distinguishes it from $\hole$. The map $n \to n_i$
is exactly the record of which binders the hole sits beneath, and it is what makes
composition well defined without $\alpha$-renaming the hole's names apart.
Consequently $\Ctx(G)$ cannot be formed in a setting where substitution is
capture-avoiding by fiat; the stage data is what replaces that discipline.
\end{remark}

\begin{remark}[Degeneration]
\label{rem:degenerate}
For Turing machines and for finite CCS without restriction, $\mathbb{F}_G$ is
terminal, every interface is a sort, and Definition~\ref{def:ctx} reduces to the
naive one. For CCS with restriction the stages are non-trivial and they matter:
$(\nu a)(\hole \mid \bar a.0)$ and $\hole \mid \bar a.0$ are distinct contexts
precisely because their holes sit at different stages. For the $\lambda$-calculus,
$n$ is the free-variable context. For $C(G)$ the adjoined sorts $\mathrm{Sig}$ and
$\mathrm{Stk}$ enlarge $S$ and leave $\mathbb{F}_G$ untouched, since signatures and
the added sorts bind nothing. No instance requires a special case.
\end{remark}

\begin{remark}[Equations on $K$ coarsen the probe]
\label{rem:Kequations}
$\Ctx(G)$ is formed modulo $\equiv$. If $K$ carries equations --- regime (b) of
Remark~\ref{rem:regimes} --- then two things happen at once, and they are the same
thing counted twice. Contexts that were positionally distinct become identified, so
$\Ctx(G)$ shrinks and the label class $L_\iota$ it induces on any target shrinks with
it; and $\smile$ is enlarged, so more introductions can meet than position licenses.
Loss of positional discrimination in the probe and leakage of ambient authority in
the theory are one phenomenon. A theory with a free $K$ is therefore a strictly
sharper probe than the same theory with equations imposed on $K$, and the two
presentations sit at different degrees.
\end{remark}

\begin{remark}[Terms are the nullary contexts]
\label{rem:nullary}
A term of sort $s$ at stage $n$ is precisely a morphism $() \to (s,n)$. Hence a
functor $\Ctx(G) \to \Ctx(H)$ encodes terms and contexts by one and the same map,
and their compatibility is functoriality rather than an axiom imposed on top. This
is the technical content of the slogan that the contexts must be mapped as well.
\end{remark}

\begin{remark}[Why a multicategory]
\label{rem:multi}
Parallel composition is genuinely binary and restriction genuinely binds across both
arguments, so the interesting contexts have several holes. A category of unary
contexts loses exactly the structure that distinguishes concurrent theories. Density
in Section~\ref{sec:conditions} is also a multi-ary condition: it asks how ambient
contexts decompose, and the decomposition is by parallel assembly.
\end{remark}

\subsection{Context-labelled transitions}
\label{sec:clts}

This is the move Section~\ref{sec:derivative} was written to prepare. Labels are not
drawn from a separately postulated alphabet of actions; they are the contexts
themselves, which is to say the locations of Definition~\ref{def:ctx}. Everything
established there about when locations separate is therefore a statement about how
much these labels can distinguish, and the reader should carry it forward:
where $\mathrm{plug}$ is injective the labelling is faithful, where it is not the
labelling collapses, and reification of the surface is what restores it.

\begin{definition}[Context-labelled transition system]
\label{def:clts}
Let $L \subseteq \Ctx_1(G)$ be a class of unary contexts containing $\hole$ and
closed under composition. The transition system $\mathrm{CLTS}(G,L)$ has terms as
states and
\[
P \xrightarrow{\;C\;} Q
\qquad\text{iff}\qquad
C \in L \text{ and } C[P] \to Q,
\]
with $\obs$ as the atomic predicate on states.
\end{definition}

\begin{definition}[$L$-bisimulation]
$\sim_L$ is the largest symmetric relation $R$ on terms such that $P \mathbin{R} Q$
implies (i) $P \obs$ iff $Q \obs$, and (ii) whenever $P \xrightarrow{C} P'$ there is
$Q \Rightarrow \xrightarrow{C} \Rightarrow Q'$ with $P' \mathbin{R} Q'$, where
$\Rightarrow$ is reflexive-transitive silent reduction.
\end{definition}

\begin{remark}[Relation to derived congruences]
\label{rem:rpo}
Taking $L = \Ctx_1(G)$ entire makes the label set unmanageably large, which is why
Leifer and Milner restrict to contexts minimal for enabling a reaction, via relative
pushouts \cite{leifer-milner,sewell,bigraphs}. We note that relativisation to the
image of a probe, introduced next, is an \emph{alternative} taming device: if the
label class is already the image of a small theory, minimality is not needed to keep
it tractable. Whether the two tamings agree on shared examples is
Problem~\ref{prob:rpo}.
\end{remark}

\subsection{Morphisms}

\begin{definition}[Ctx-morphism]
\label{def:ctxmor}
A \emph{Ctx-morphism} $\iota : G \to H$ is a symmetric pseudofunctor of
multicategories $\iota : \Ctx(G) \to \Ctx(H)$ such that
\begin{enumerate}[label=(\roman*),leftmargin=2.2em]
\item $\iota(\hole) \bisim \hole$ and
$\iota\bigl(C \circ (D_1,\dots,D_n)\bigr) \bisim \iota(C) \circ (\iota D_1,\dots,\iota D_n)$,
the comparisons being the structure 2-cells of the pseudofunctor;
\item $\iota$ reflects and preserves the primitive observation:
$C[\vec P] \obs$ iff $\iota(C)[\iota \vec P] \obs$;
\item $\iota$ is \emph{stage-indexed}: it is accompanied by a functor
$\varphi_\iota : \mathbb{F}_G \to \mathbb{F}_H$ over which it acts on interfaces,
$(s,n) \mapsto (\iota s, \varphi_\iota n)$, compatibly with the maps $n \to n_i$
recorded at holes.
\end{enumerate}
Write $\CtxGSLT$ for the resulting bicategory, with 2-cells the bisimulations
between images.
\end{definition}

Note that ``pseudo'' is not fussiness. Encodings routinely preserve context
composition only up to administrative reduction --- Milner's forwarders are exactly
that --- so strict functoriality would exclude the examples.

\begin{definition}[Induced label class and relative bisimulation]
\label{def:relbis}
For $\iota : G \to H$ put $L_\iota := \iota\bigl(\Ctx_1(G)\bigr)$, closed under
composition. Two $H$-terms $X,Y$ are \emph{$\iota$-indistinguishable}, written
$X \sim_\iota Y$, if they are bisimilar in $\mathrm{CLTS}(H, L_\iota)$.
\end{definition}

This is the formal content of the observation that when $\pi$ is encoded into rho,
the encoding may be probed only by the images of $\pi$ contexts.

\begin{proposition}[Preservation is now structural]
\label{prop:preserve}
For any Ctx-morphism $\iota$, the assignment $P \mapsto \iota(P)$ is a morphism of
transition systems
$\mathrm{CLTS}(G,\Ctx_1(G)) \to \mathrm{CLTS}(H,L_\iota)$.
Consequently $P \sim_{\Ctx_1(G)} Q$ implies $\iota P \sim_\iota \iota Q$.
\end{proposition}

\begin{proof}
If $P \xrightarrow{C} Q$ then $C[P] \to Q$; applying $\iota$ and using (i) of
Definition~\ref{def:ctxmor}, $\iota(C)[\iota P]$ is bisimilar to $\iota(C[P])$, which
by (ii) and the fact that pseudofunctors preserve the reduction witnessed by $\obs$
reduces to a term bisimilar to $\iota Q$. Labels are transported because
$\iota(C) \in L_\iota$ by definition. The consequence is the standard transport of
bisimilarity along a morphism of transition systems.
\end{proof}

\begin{remark}
\label{rem:kernel-congruence}
Proposition~\ref{prop:preserve} costs nothing because functoriality was built into
Definition~\ref{def:ctxmor}. This is the first dividend of taking contexts as
primary: the compositionality condition that had to be imposed as an axiom in the
naive setting is now definitional, and the induced kernel (Definition~\ref{def:F})
is automatically a congruence.
\end{remark}

\subsection{Adequacy of the probe}

Relativising the label class introduces a new way to cheat: the fewer contexts one
admits, the coarser the induced equivalence and the easier preservation becomes. In
the limit, admitting no contexts makes everything equivalent. The invariant of
interest is therefore always the \emph{pair} $(\iota, L_\iota)$, never the encoding
alone, and the probe must be pinned down independently.

\begin{definition}[Self-separating probe]
\label{def:selfsep}
$G$ is \emph{self-separating} if $\sim_{\Ctx_1(G)}$ coincides with the intended
weak bisimulation $\bisim_G$.
\end{definition}

If a probe is not self-separating then a statement of the form ``the encoding
preserves bisimulation over encoded $G$ contexts'' preserves something other than
what $G$ means by equivalence, and the whole exercise is void. Self-separation is a
hypothesis on the probe alone, verified once per probe, and is the standard question
of whether a derived context congruence agrees with barbed congruence
\cite{leifer-milner,sewell}.

\section{The two conditions}
\label{sec:conditions}

\subsection{Hosting}

\begin{definition}[Condition (F): hosting]
\label{def:F}
A Ctx-morphism $\iota : G \to H$ \emph{hosts} $G$, written $\iota \models (F)$, if it
is faithful up to bisimulation: for all parallel $C, D$ in $\Ctx(G)$,
\[
\iota(C) \sim_\iota \iota(D) \implies C \sim_{\Ctx_1(G)} D .
\]
The \emph{kernel} of $\iota$ is the congruence
$C \approx_\iota D \iff \iota(C) \sim_\iota \iota(D)$; condition (F) says the kernel
is as fine as possible.
\end{definition}

\begin{proposition}[Non-degeneracy]
\label{prop:nondeg}
The constant assignment $\iota(-) = 0$ does not satisfy (F) whenever $G$ has two
non-bisimilar terms.
\end{proposition}

\begin{proof}
All nullary contexts have the same image, so the kernel is the total relation,
which is not $\sim_{\Ctx_1(G)}$ under the hypothesis.
\end{proof}

So the degeneracy that motivated the search is excluded by (F) alone, and by a
condition that mentions neither syntax nor logic.

\subsection{Exhausting}

Hosting says the target is at least as expressive as the probe. The complementary
question --- whether the target is anything \emph{more} --- is density.

\begin{definition}[Restricted nerve]
\label{def:nerve}
For $\iota : G \to H$ define
\[
\Nerve_\iota : \Ctx(H) \longrightarrow \Psh\bigl(\Ctx(G)\bigr),
\qquad
\Nerve_\iota(X)(C) = \Sim\bigl(\iota C,\, X\bigr),
\]
where $\Sim(-,-)$ is the set of simulations, taken up to $\sim_\iota$.
\end{definition}

\begin{definition}[Condition (D): exhausting]
\label{def:D}
$\iota \models (D)$ if $\iota$ is \emph{dense}, i.e.\ $\Nerve_\iota$ is fully
faithful. Equivalently, every $H$-context is a canonical colimit of the diagram of
$\iota$-images mapping into it.
\end{definition}

\begin{proposition}[Operational reading]
\label{prop:Dop}
$\iota \models (D)$ iff every $H$-context $D$ is $\sim_\iota$-indistinguishable from
a context assembled out of $\iota$-images: there is a finite diagram
$C_1,\dots,C_k$ in $\Ctx(G)$ and an assembly of $\iota C_1, \dots, \iota C_k$ by
parallel composition, restriction, and hole-filling that is $\sim_\iota$-equivalent
to $D$.
\end{proposition}

\begin{design}[Which colimits]
\label{des:colim}
Proposition~\ref{prop:Dop} is stated for the class of \emph{assembly colimits}:
those generated by parallel composition, restriction, and hole-filling, closed under
the reduction-directed joins needed to make the presentation stable under $\to$.
Pinning this class down precisely is the main piece of unfinished business in this
note; see Problem~\ref{prob:colim}. We use the operational reading as the working
definition throughout and flag every place where the choice matters.
\end{design}

\begin{proposition}[Composition]
\label{prop:compose}
Condition (F) is stable under composition. Condition (D) is stable under composition
of morphisms preserving assembly colimits.
\end{proposition}

\begin{proof}
For (F), faithfulness composes. For (D), if every $H$-context is an assembly of
$\iota$-images and every $K$-context is an assembly of $\kappa$-images, then every
$K$-context is an assembly of assemblies of $\kappa\iota$-images, and an assembly of
assemblies is an assembly provided $\kappa$ carries the inner assemblies to
assemblies --- which is the stated hypothesis.
\end{proof}

\begin{remark}
The hypothesis in Proposition~\ref{prop:compose} is not idle: density is famously not
stable under composition in general, and the failure is exactly failure of
cocontinuity. This is a real constraint on the theory, not a technicality to be
assumed away.
\end{remark}

\subsection{Relative universality}

\begin{definition}[The probe order]
\label{def:order}
Write $G \preceq H$ if some Ctx-morphism $\iota : G \to H$ satisfies (F); say $H$ is
\emph{$G$-universal}. Write $G \equiv H$ if some $\iota$ satisfies both (F) and (D);
say $H$ is \emph{$G$-exact}. Write $G \prec H$ if $G \preceq H$ but no $\iota$
satisfying (F) also satisfies (D).
\end{definition}

\begin{proposition}
$\preceq$ is a preorder on theories, and $\equiv$ is the induced equivalence on
morphisms preserving assembly colimits.
\end{proposition}

\begin{proof}
Reflexivity by the identity, transitivity by Proposition~\ref{prop:compose}.
\end{proof}

So theories carry \emph{degrees} under $\preceq$, and $\prec$ is a strictness
relation with a positive criterion: hosted but not exhausted. The observation that
arbitrary rho contexts see strictly more than $\pi$ contexts is now a
theorem-shaped assertion --- that no (F)-morphism $\pi \to \text{rho}$ is dense ---
rather than an obstruction to the existence of morphisms.

\subsection{Image factorisation}

Condition (F) is too strong to demand of every encoding, since Milner's fails it.
The general statement is a factorisation.

\begin{proposition}[Image factorisation]
\label{prop:factor}
Every Ctx-morphism $\iota : G \to H$ factors as
\[
\begin{tikzcd}[column sep=2.6em]
G \arrow[r, two heads, "q"] & G/{\approx_\iota} \arrow[r, hook, "\bar\iota"] & H
\end{tikzcd}
\]
where $q$ is the quotient by the kernel and $\bar\iota \models (F)$.
\end{proposition}

\begin{proof}
The kernel $\approx_\iota$ is a congruence by
Remark~\ref{rem:kernel-congruence}, so the quotient multicategory exists; $\iota$
factors through it by construction, and the induced $\bar\iota$ is faithful by
definition of the kernel.
\end{proof}

\begin{example}[Milner, located]
\label{ex:milner}
For Milner's call-by-name encoding, $q$ is the quotient of $\Lambda$ by L\'evy--Longo
tree equality and $\bar\iota$ hosts the quotient inside $\pi$. The encoding is thus
completely described by the pair: a nontrivial quotient recording under-discrimination,
followed by a hosting embedding. Neither half is visible in
Definition~\ref{def:naive}.
\end{example}

\begin{remark}[The full invariant]
An encoding is characterised by three data: the kernel (how much of the source is
lost), the failure of density (how much of the target is unseen), and --- once
Section~\ref{sec:meta-grade} is in play --- the overhead. Reporting only that ``an
encoding exists'' discards all three.
\end{remark}

\section{Refining Turing completeness}
\label{sec:tc}

\subsection{The classical notion as a trace collapse}

\begin{definition}[Trace collapse]
\label{def:collapse}
The \emph{trace collapse} of $\CtxGSLT$ is the bicategory obtained by replacing the
2-cells --- bisimulations --- by trace inclusions, so that two terms are identified
when they admit the same sequences of labels.
\end{definition}

By Remark~\ref{rem:machines} the Turing machine is not an object of the category, so
before it can serve as a probe at all it must be \emph{presented} as one.

\begin{definition}[A presentation of $\mathrm{TM}$]
\label{def:tmpres}
A \emph{process presentation} of Turing machines is an interactive GSLT
$\mathrm{TM}^\sharp$ whose interacting sort contains both automaton and tape --- for
instance by encoding the tape as a process --- with $K$ the constructor placing an
automaton beside a tape, and whose interaction rule is the machine's transition
relation.
\end{definition}

Such a presentation exists, and that is the whole difficulty: it is a choice.
Different encodings of the tape yield different $\Ctx(\mathrm{TM}^\sharp)$ and hence
different probes, and nothing internal to the machine model selects among them. The
classical yardstick is therefore not merely coarse; it is not even well defined as a
probe until an extraneous decision has been made.

\begin{proposition}[Recovering Turing completeness]
\label{prop:tc}
Fix a process presentation $\mathrm{TM}^\sharp$. In the trace collapse, $H$ is
$\mathrm{TM}^\sharp$-universal if and only if $H$ is Turing complete in the classical
sense, and the criterion is independent of which presentation was fixed.
\end{proposition}

\begin{proof}[Proof sketch]
In the trace collapse the branching structure of $\Ctx(\mathrm{TM}^\sharp)$ is
invisible, so the only surviving content of faithfulness on the nullary component is
that distinct input/output functions have distinct images --- i.e.\ that $H$ realises
every partial recursive function. Conversely a universal machine simulated in $H$
supplies the morphism. Independence of the presentation holds because any two process
presentations of the same machine model compute the same functions, and at trace level
that is all that is being compared.
\end{proof}

\begin{remark}[Presentation-dependence is the diagnostic]
\label{rem:presdep}
The independence clause of Proposition~\ref{prop:tc} holds \emph{only} in the trace
collapse. Before collapsing, distinct process presentations of the same machine model
sit at distinct degrees, since they induce distinct context multicategories. The
classical notion is thus well defined precisely to the extent that it is blind, and
the coincidence of those two facts is not an accident: what makes it insensitive to
the choice of presentation is exactly what makes it insensitive to branching.
\end{remark}

\begin{remark}[Why the classical notion is blind]
\label{rem:blind}
Above the Turing threshold the trace-collapsed order has a single degree: any two
Turing-complete theories host one another at trace level. The branching-time order
does not collapse, and the separation results \cite{palamidessi,gorla,peters-glabbeek}
are exactly exhibitions of that non-collapse. Classical Turing completeness is thus
the shadow of $\preceq$ under Definition~\ref{def:collapse}, and using it as a
hypothesis for a bisimulation-level conclusion is using the shadow to constrain the
object.
\end{remark}

\subsection{The ladder}

The following stratification replaces the single predicate. Each level adds one
requirement; the classical notion is level 0, and level 4 is what a genuine claim of
behavioural universality amounts to.

\begin{center}
\begin{tabular}{@{}clp{6.4cm}@{}}
\toprule
Level & Condition & Comment \\
\midrule
0 & functional universality & classical Turing completeness; trace level \\
1 & sound simulation & $\iota$ is a Ctx-morphism; no faithfulness claimed \\
2 & hosting (F) & kernel is trivial; source runs faithfully in target \\
3 & internal representability (I) & the encoding is realised by a single target
term applied to codes, not by a metatheoretic family \\
4 & exhausting (D) & target adds nothing the source cannot assemble \\
\bottomrule
\end{tabular}
\end{center}

\begin{definition}[Condition (I): internal representability]
\label{def:I}
$\iota \models (I)$ if there is a coding $\code{-}$ of $G$-terms as $H$-data,
definable in $H$, and a single $H$-term $U_\iota$ with
$U_\iota\langle \code{P}\rangle \sim_\iota \iota(P)$ for every $G$-term $P$.
\end{definition}

Condition (I) is what turns a translation into an interpreter, and it factors
exactly along the ladder of Section~\ref{sec:naive}.

\begin{proposition}[Continuity and reflection give (I)]
\label{prop:contrefl}
Let $G$ and $H$ be continued (Definition~\ref{def:cigslt}) and let $H$ be
reflective, in the sense that it carries a term-level quoting operator $@$ with
$*@P \equiv P$. Then any Ctx-morphism $\iota : G \to H$ satisfies (I).
\end{proposition}

\begin{proof}[Proof sketch]
By clause (ii) and Proposition~\ref{prop:effctx}, $\Ctx(G)$ admits codes, so $\code{-}$ is
a computable map from $G$-terms to finite data. By reflection of $H$, that data is
representable as an $H$-name, and $*$ recovers the process it names; $U_\iota$ is
then the $H$-term that decodes and dispatches, definable because the guards it must
decide are decidable by clause (ii) of $H$.
\end{proof}

The two hypotheses do distinct work, and separating them is the point: continuity
of $G$ supplies canonical representatives to be coded, reflection of $H$ supplies
an internal home for the codes. In $\pi$ the second fails --- processes are not
transmissible data, so any coding is metatheoretic --- while in rho both hold by
construction. Non-naturality of an internalisation family and externality of its
coding are the same defect seen twice, and Proposition~\ref{prop:contrefl} says
exactly which hypothesis repairs it.

\subsection{Separations}

\begin{center}
\begin{tabular}{@{}lcccl@{}}
\toprule
Probe / target & (F) & (I) & (D) & Obstruction \\
\midrule
$\mathrm{TM}^\sharp \to$ CCS+rec & \checkmark & \checkmark & --- & no name passing; probe needs a presentation \\
$\pi \to$ CCS+rec & --- & --- & --- & fixed interface \\
$\lambda \to \pi$ & --- & \checkmark & --- & kernel is L\'evy--Longo \\
$\pi \to$ rho & \checkmark & \checkmark & --- & $@$ not assemblable \\
rho $\to$ rho & \checkmark & \checkmark & \checkmark & --- \\
\bottomrule
\end{tabular}
\end{center}

\subsection{Separation at the level of stages}

The entry for $\pi \to \mathrm{CCS}$ deserves its own statement, since it is the
cleanest separation the order provides and it is invisible classically. With
Definition~\ref{def:stage} in hand it can be argued at the level of stage
categories, without descending to a particular pair of terms.

\begin{definition}[Stage-changing transition]
\label{def:stagechange}
A transition $P \xrightarrow{C} Q$ is \emph{stage-changing} if the stage of $Q$
properly extends that of $P$ in $\mathbb{F}_G$.
\end{definition}

Bound output in $\pi$ is stage-changing: extruding a private name enlarges the
observer's stage. Restriction in CCS is not: $(\nu a)P$ never emits $a$, so every
CCS transition is stage-preserving, and $\mathbb{F}_{\mathrm{CCS}} = \mathbb{I}$
contains only injections.

\begin{proposition}[Hosting constrains the stage functor]
\label{prop:stagefunctor}
If $\iota : G \to H$ satisfies (F), then $\varphi_\iota : \mathbb{F}_G \to
\mathbb{F}_H$ is faithful and preserves the substitution action: for every map
$\sigma$ in $\mathbb{F}_G$ whose induced action separates two $G$-terms,
$\varphi_\iota \sigma$ separates their images.
\end{proposition}

\begin{proof}[Proof sketch]
Suppose $\varphi_\iota \sigma = \varphi_\iota \tau$ for distinct $\sigma, \tau$
separating $P$ from $Q$. The $G$-contexts realising $\sigma$ and $\tau$ then have
$\sim_\iota$-equal images, so the kernel of $\iota$ identifies them, contradicting
(F). Preservation of the action is the same argument applied to the nullary
component.
\end{proof}

\begin{corollary}[$\pi \not\preceq \mathrm{CCS}$]
\label{cor:pinotccs}
No Ctx-morphism $\pi \to \mathrm{CCS}{+}\mathrm{rec}$ satisfies (F), although
$\mathrm{CCS}{+}\mathrm{rec}$ is Turing complete.
\end{corollary}

\begin{proof}[Proof sketch]
$\mathbb{F}_\pi$ contains non-injective maps: a $\pi$ context may deliver a name
that identifies two names previously distinct, turning an inert composition into a
redex. Every map of $\mathbb{F}_{\mathrm{CCS}} = \mathbb{I}$ is injective, so
$\varphi_\iota$ must send a non-monic to a monic and therefore identifies maps that
$\mathbb{F}_\pi$ distinguishes. By Proposition~\ref{prop:stagefunctor} this
contradicts (F).
\end{proof}

\begin{remark}[Residual hypothesis]
\label{rem:residual}
The sketch rules out encodings that are stage-indexed in the sense of
Definition~\ref{def:ctxmor}(iii). It does not by itself rule out an encoding that
pre-allocates a fixed finite pool of CCS channels and simulates identification by
case analysis over the pool. Excluding those requires the additional hypothesis
that $\varphi_\iota$ be finitary and uniform --- essentially the finite-support
condition of nominal sets. We record this as a hypothesis rather than pretending it
is discharged; it is the same gap that the uniformity criteria of \cite{gorla} are
designed to close, and it is the one place where this note's machinery does not yet
subsume the checklist.
\end{remark}

Behavioural universality therefore forces name passing; under (I) it forces the
coding to land in the name sort, which is reflection.

\section{Two meta-circular interpreters}
\label{sec:meta}

We come to the structural payoff. Conditions (F) and (D) look like two unrelated
technical hygiene requirements. They are in fact the two halves of mutual
meta-circularity: one gives an interpreter running the source inside the target, the
other an interpreter running the target inside the source.

\subsection{The source, inside the target}

\begin{theorem}[Hosting yields an interpreter for the source]
\label{thm:src-in-tgt}
Let $\iota : G \to H$ satisfy (F) and (I), and let $G$ be self-separating. Then
$U_\iota$ is an interpreter for $G$ written in $H$, correct up to $\sim_\iota$: for
every $G$-term $P$ and every $G$-context $C$,
\[
U_\iota\langle \code{C[P]}\rangle \;\sim_\iota\; \iota(C)\bigl[\,U_\iota\langle\code{P}\rangle\,\bigr],
\]
and $P \bisim_G Q$ iff
$U_\iota\langle\code{P}\rangle \sim_\iota U_\iota\langle\code{Q}\rangle$.
\end{theorem}

\begin{proof}
The displayed equation is Definition~\ref{def:I} composed with functoriality (i) of
Definition~\ref{def:ctxmor}. The right-to-left direction of the second claim is
Proposition~\ref{prop:preserve}; the left-to-right direction is (F) together with
self-separation of $G$, which identifies $\sim_{\Ctx_1(G)}$ with $\bisim_G$.
\end{proof}

Write $E_{G\to H} := U_\iota$ for this interpreter.

\subsection{The target, inside the source}

\begin{theorem}[Exhausting yields an interpreter for the target]
\label{thm:tgt-in-src}
Let $\iota : G \to H$ satisfy (D) with $G$ and $H$ continued. Then the assembly
presentations of Proposition~\ref{prop:Dop} are effective, and there is a $G$-term
$V_\iota$ interpreting $H$ inside $G$: for every $H$-context $D$ and $H$-term $X$,
\[
V_\iota\langle \code{D},\code{X}\rangle \;\bisim_G\; \text{the }G\text{-context assembly } \delta(D) \text{ applied to } \delta(X),
\]
where $\delta$ is the computable procedure taking a code for an $H$-context to a
finite assembly of $\iota$-images $\sim_\iota$-equivalent to it.
\end{theorem}

\begin{proof}[Proof sketch]
By (D), each $H$-context is $\sim_\iota$-equivalent to an assembly of $\iota$-images.
By clause (ii) for $H$ and Proposition~\ref{prop:effctx}, $H$-contexts admit codes, and by
clause (ii) for $G$ the search for an assembly ranges over an effectively enumerable set of
$G$-context representatives with decidable equality, so $\delta$ is computable; by faithfulness of $\iota$
on the assembled pieces --- which is the fully faithful half of
Definition~\ref{def:D} --- the $G$-contexts underlying those images are determined up
to $\sim_{\Ctx_1(G)}$. $V_\iota$ therefore decodes, recovers the underlying
$G$-contexts, runs them under $G$'s own dynamics, and reassembles. Correctness up to
$\bisim_G$ follows from self-separation of $G$ and Proposition~\ref{prop:preserve}
applied in the reverse direction along the fully faithful nerve.
\end{proof}

Write $E_{H\to G} := V_\iota$.

\begin{remark}[Where the asymmetry lives]
Theorem~\ref{thm:src-in-tgt} needs (I) and Theorem~\ref{thm:tgt-in-src} needs
effectiveness of $\delta$. Both are ``the encoding is not merely a metatheoretic
function'' conditions, imposed once on each side, and both are now consequences
rather than hypotheses: (I) follows from continuity together with reflection of the
target (Proposition~\ref{prop:contrefl}), and effectiveness of $\delta$ follows from
the sections of both theories (Proposition~\ref{prop:effctx}). This is the whole
purpose of rung three of Section~\ref{sec:naive}. In a reflective continued theory
the same operator supplies both.
\end{remark}

\subsection{The two round trips}

\begin{corollary}[Two meta-circular interpreters]
\label{cor:meta}
Suppose $G$ and $H$ are continued and self-separating, $H$ is reflective, and
$\iota : G \to H$ satisfies (F) and (D). Then (I) and effective assembly hold by
Propositions~\ref{prop:contrefl} and \ref{prop:effctx}, and
\[
M_G \;:=\; E_{H\to G} \circ E_{G\to H} \;:\; G \rightsquigarrow G,
\qquad
M_H \;:=\; E_{G\to H} \circ E_{H\to G} \;:\; H \rightsquigarrow H
\]
are interpreters for $G$ in $G$ and for $H$ in $H$ respectively, satisfying
\[
M_G \;\bisim_G\; \mathrm{id}_G,
\qquad
M_H \;\sim_\iota\; \mathrm{id}_H .
\]
$M_G$ is the source looping through the target; $M_H$ is the target looping through
the source.
\end{corollary}

\begin{proof}
Composition of the two theorems. $M_G$ takes a $G$-term to its code, runs it under
$H$'s dynamics as simulated by $V_\iota$, and returns; each leg preserves and
reflects the relevant equivalence, so the composite is behaviourally the identity.
Symmetrically for $M_H$, with $\sim_\iota$ in place of $\bisim_H$ because $H$ is
observed only through $L_\iota$.
\end{proof}

\begin{definition}[Meta-circularity]
\label{def:metacirc}
An endo-interpreter $M : G \rightsquigarrow G$ is \emph{meta-circular} if it is
syntactically non-trivial --- $M(P) \not\equiv P$ for some $P$ --- and behaviourally
the identity, $M \bisim \mathrm{id}$.
\end{definition}

Corollary~\ref{cor:meta} says that exactness of a probe relation manufactures a
meta-circular interpreter on each side. This is the precise sense in which the two
conditions belong together: (F) alone gives a one-way hosting with no return
journey, (D) alone gives a decomposition with nothing to decompose faithfully, and
only the conjunction closes the loop.

\begin{remark}[Monad and comonad]
\label{rem:monad}
$E_{G\to H} \dashv E_{H\to G}$ is an adjunction in the simulation bicategory whose
induced monad on $G$ is $M_G$ and whose induced comonad on $H$ is $M_H$. Both are
idempotent up to bisimulation, so the adjunction is an adjoint equivalence. In
such situations the degeneracy is usually noted with unease --- if both unit and
counit are invertible, the word ``adjunction'' does no work. The reading offered here
is the opposite: the degeneracy \emph{is} the content. An adjoint equivalence between
two theories in the simulation bicategory is precisely the statement that each can
meta-circularly interpret the other. What is missing in such cases is not
non-degeneracy but a grading, supplied next.
\end{remark}

\subsection{Grading: meta-circularity is not free}
\label{sec:meta-grade}

\begin{definition}[Graded interpreter]
\label{def:graded}
$E$ has \emph{grade} $g : \mathbb{N} \to \mathbb{N}$ if every labelled step of the
source is matched by at most $g(n)$ steps of the target, $n$ the size of the
argument. Write $G \preceq_g H$ for hosting with grade $g$.
\end{definition}

\begin{proposition}[The cost of the loop]
\label{prop:loop-cost}
If $E_{G\to H}$ has grade $g$ and $E_{H\to G}$ has grade $h$, then $M_G$ has grade
$g \circ h$. Hence $M_G \bisim \mathrm{id}_G$ but $M_G \not\bisim_{1} \mathrm{id}_G$
whenever $g\circ h$ is not the identity: the round trip is behaviourally invisible and
resource-visible.
\end{proposition}

\begin{corollary}[Graded lax coreflection]
\label{cor:lax}
In the graded simulation bicategory, $E_{G\to H} \dashv E_{H\to G}$ is a lax
coreflection: the unit has grade $1$ and the counit has grade $g \circ h$. The
failure to be an equivalence is exactly the interpretive overhead.
\end{corollary}

Ungraded, the adjunction collapses and says nothing; graded, its counit is a
measurement, and the thing being measured is the price of self-interpretation. The
grading is stated here because it costs nothing to state and because the interpretive
overhead it names is the only quantity in this note with a plausible claim to being
measured rather than merely compared; nothing else below depends on it.

\section{Open problems}
\label{sec:open}

\begin{problem}[Assembly colimits]
\label{prob:colim}
Pin down the class of colimits in Design~\ref{des:colim}. Density is only as sharp as
the colimits one admits, and admitting too many trivialises (D) while admitting too
few makes it unsatisfiable. The candidate is the class generated by parallel
composition, restriction, and hole-filling, closed under reduction-directed joins;
what must be checked is that this class is stable under Ctx-morphisms preserving
$\obs$.
\end{problem}

\begin{problem}[Self-separation of $\pi$]
\label{prob:rpo}
Is $\pi$ self-separating in the sense of Definition~\ref{def:selfsep}? Equivalently,
does the context-labelled bisimulation over all unary $\pi$ contexts agree with
barbed congruence? If the RPO-derived labels of \cite{leifer-milner} suffice, does the
relativised label class $L_\iota$ agree with them on shared examples?
\end{problem}

\begin{problem}[The rho probe of $\lambda$]
Compute the kernel of the canonical encoding $\lambda \to \text{rho}$. Since
reflection permits the tree structure of a term to be internalised, the kernel may be
strictly finer than L\'evy--Longo. If it is $\beta$-conversion itself, that is a
sharp statement of what reflection adds over name passing.
\end{problem}

\begin{problem}[$\pi \prec$ rho]
\label{prob:pirho}
Exhibit a rho context --- $@$, or the nearness operator of
Definition~\ref{def:surface} --- and
prove that it is not $\sim_\iota$-equivalent to any assembly of encoded $\pi$
contexts. This is the first genuine strictness result the order should deliver, and
it is the one the classical yardstick is constitutionally unable to see.

Remark~\ref{rem:stages} suggests attacking it at the level of stages rather than
terms, as Corollary~\ref{cor:pinotccs} does for the CCS case: $\mathbb{F}_\pi$ is a
category of finite sets, whereas $\mathbb{F}_{\mathrm{rho}}$ satisfies
$\mathcal{N} \cong @\,\mathrm{Proc}$ and is therefore a fixed point no functor from
finite sets can reach. If the argument goes through, both separations become
statements about binding signatures alone, and the probe order acquires a
computable lower bound.
\end{problem}

\begin{problem}[Discharging Remark~\ref{rem:residual}]
Formulate the finitary uniformity hypothesis on $\varphi_\iota$ precisely --- most
likely as finite support in the sense of nominal sets --- and verify that it holds
for every Ctx-morphism arising from an internally representable encoding in the
sense of Definition~\ref{def:I}. If (I) implies it, the residual gap closes and the
separation arguments become unconditional.
\end{problem}

\begin{problem}[Well-foundedness of the graded order]
Is $\preceq_g$ well-founded? A descending chain of theories each hosting the next
with strictly better grade would be a hierarchy of increasingly efficient
self-interpretation, and its existence or absence is the quantitative analogue of the
degree structure.
\end{problem}

\begin{problem}[Probe profiles]
A theory's invariant under this analysis is not a single degree but the profile of
its kernels and densities across a family of probes. Is the profile of rho determined
by finitely many probes?
\end{problem}

\begin{thebibliography}{99}

\bibitem{cost}
L.\ G.\ Meredith.
\newblock \emph{Continued Interactive GSLTs and the Cost Endofunctor: a construction
one level up from cost-accounted Rholang}.
\newblock F1R3FLY.io, 2026.

\bibitem{rho}
L.\ G.\ Meredith and M.\ Radestock.
\newblock A reflective higher-order calculus.
\newblock \emph{Electronic Notes in Theoretical Computer Science}, 2005.

\bibitem{milner-fap}
R.\ Milner.
\newblock Functions as processes.
\newblock \emph{Mathematical Structures in Computer Science}, 1992.

\bibitem{sangiorgi-lazy}
D.\ Sangiorgi.
\newblock The lazy lambda calculus in a concurrency scenario.
\newblock \emph{Information and Computation}, 1994.

\bibitem{sw}
D.\ Sangiorgi and D.\ Walker.
\newblock \emph{The $\pi$-Calculus: A Theory of Mobile Processes}.
\newblock Cambridge University Press, 2001.

\bibitem{leifer-milner}
J.\ Leifer and R.\ Milner.
\newblock Deriving bisimulation congruences for reactive systems.
\newblock In \emph{CONCUR}, 2000.

\bibitem{sewell}
P.\ Sewell.
\newblock From rewrite rules to bisimulation congruences.
\newblock \emph{Theoretical Computer Science}, 2002.

\bibitem{bigraphs}
R.\ Milner.
\newblock \emph{The Space and Motion of Communicating Agents}.
\newblock Cambridge University Press, 2009.

\bibitem{palamidessi}
C.\ Palamidessi.
\newblock Comparing the expressive power of the synchronous and the asynchronous
$\pi$-calculi.
\newblock \emph{Mathematical Structures in Computer Science}, 2003.

\bibitem{gorla}
D.\ Gorla.
\newblock Towards a unified approach to encodability and separation results for
process calculi.
\newblock \emph{Information and Computation}, 2010.

\bibitem{peters-glabbeek}
K.\ Peters and R.\ J.\ van Glabbeek.
\newblock Analysing and comparing encodability criteria for process calculi.
\newblock In \emph{EXPRESS/SOS}, 2015.

\bibitem{mcbride}
C.\ McBride.
\newblock The derivative of a regular type is its type of one-hole contexts.
\newblock Unpublished manuscript, 2001.

\bibitem{huet}
G.\ Huet.
\newblock The zipper.
\newblock \emph{Journal of Functional Programming}, 1997.

\bibitem{kelly}
G.\ M.\ Kelly.
\newblock \emph{Basic Concepts of Enriched Category Theory}.
\newblock Cambridge University Press, 1982.

\bibitem{cwm}
S.\ Mac Lane.
\newblock \emph{Categories for the Working Mathematician}.
\newblock Springer, 2nd edition, 1998.

\end{thebibliography}

\end{document}
